\begin{table}[htbp]
\centering
\caption{Standardized Effect Sizes}
\label{tab:sde}
\begin{threeparttable}
\begin{tabular}{lS[table-format=-1.2]S[table-format=1.2]S[table-format=1.2]S[table-format=-1.3]S[table-format=1.3]l}
\toprule
Outcome & {$\hat{\beta}$} & {SE} & {SD($Y$)} & {SDE} & {SE(SDE)} & Classification \\
\midrule
\multicolumn{7}{l}{\textit{Panel A: Pooled}} \\
Deceased donors pmp & -1.63 & 1.85 & 3.74 & -0.436 & 0.494 & Large negative \\
Transplants pmp & -15.10 & 6.19 & 11.22 & -1.346 & 0.552 & Large negative \\
Living donors pmp & -3.32 & 3.90 & 8.82 & -0.377 & 0.442 & Large negative \\
\midrule
\multicolumn{7}{l}{\textit{Panel B: Heterogeneous (Early vs.\ Late Adopters)}} \\
Early adopters (W, E) & -0.16 & 2.11 & 1.96 & -0.082 & 1.077 & Moderate negative \\
Late adopters (S, NI) & -3.89 & 1.54 & 4.29 & -0.907 & 0.360 & Large negative \\
\bottomrule
\end{tabular}
\begin{tablenotes}[flushleft]\footnotesize
\item \textit{Notes:} \textbf{Country:} United Kingdom (England, Wales, Scotland, Northern Ireland). \textbf{Research question:} Does opt-out (deemed consent) organ donation legislation increase deceased donor rates and transplantation? \textbf{Policy mechanism:} Deemed consent laws shift the legal default from requiring explicit opt-in registration to presuming consent unless the individual has actively opted out, effectively reclassifying non-registrants as potential donors and transferring the decision burden to families at the point of death. \textbf{Outcome definition:} Deceased donors per million population (dd\_pmp), counting all donors from whom at least one organ was retrieved for transplant, per NHSBT definition. \textbf{Treatment:} Binary indicator equal to one from the first full financial year after deemed consent legislation took effect in each nation (Wales 2016/17, England 2020/21, Scotland 2021/22, Northern Ireland 2023/24). \textbf{Data:} NHSBT Annual Activity Reports, financial years 2015/16--2024/25, nation-year panel, 40 observations (4 nations $\times$ 10 years). \textbf{Method:} Two-way fixed effects (nation + year FE), heteroskedasticity-robust SEs, randomization inference (5,000 permutations of treatment assignment). \textbf{Sample:} All four UK nations; treatment is staggered across nations (Wales first, NI last). SDE $= \hat{\beta} / \text{SD}(Y)$ where SD($Y$) is the pre-treatment standard deviation. Classification refers to magnitude, not statistical significance: Large ($|$SDE$| > 0.15$), Moderate ($0.05$--$0.15$), Small ($0.005$--$0.05$), Null ($< 0.005$).
\end{tablenotes}
\end{threeparttable}
\end{table}
